% generated by roam.py handbook export -- не редактировать, править заметки
% 20261007202246-алгоритм_backtracking_для_подбора_шага_армихо_вульфа.org (body 4acd14827128)

\begin{method}[{Алгоритм Backtracking для подбора шага Армихо-Вульфа}]\label[method]{note:F16921E1-67C5-4EDE-A55F-AD11D8257A69}
  Алгоритм ищет шаг, который удовлетворяет \hyperref[note:15611E05-13FC-4FB0-AE00-A4BB857A29E6]{условию Армихо}.
\end{method}

\paragraph*{Алгоритм}
\begin{Verbatim}[fontsize=\small]
# Условные обозначения:
#   phi(a) = f(x + a*d) - функция по направлению d
#   dphi(a) = grad_f(x+a*d) @ d - производная по направлению d

def backtrack(phi, dphi0, alpha0=1.0, c1=1e-4, beta=0.5):
    """Шаг с условием Армихо. dphi0 = phi'(0) < 0, 0 < beta < 1"""
    phi0 = phi(0.0)
    alpha = alpha0
    while phi(alpha) > phi0 + c1 * alpha * dphi0: # Условие армихо нарушено
        alpha *= beta
    return alpha
\end{Verbatim}

\paragraph*{Почему алгоритм остановится?}
В 1 шаге доказательства теоремы \hyperref[note:C143295F-C422-4FEE-95AD-3D6B620BC0E8]{Существование шага для Армихо-Вульфа} мы доказали, что условие Армихо выполнено во всех точках интервала $(0;\delta)$. А так как $\beta^k\alpha_0 \rightarrow 0$, то при достаточно большом $k$, оно окажется в нужном интервале.

\paragraph*{Почему шаг достаточно большой?}
По \hyperref[note:FFEA1F1D-B70C-4B4A-9386-8344A2A3BBC4]{оценке убывания за шаг} для $y = x - \alpha \nabla f(x)$, где $\alpha$~--- шаг, который мы получили:
\[
f(x - \alpha \nabla f(x)) \le f(x) - \alpha \left( 1 - \frac{L\alpha}{2} \right) \|\nabla f(x)\|^2
\]

Зададим $\phi(\alpha) = f(x - \alpha \nabla f(x))$. Тогда \hyperref[note:15611E05-13FC-4FB0-AE00-A4BB857A29E6]{условие Армихо} для $\alpha_0 = 0$, $\alpha_k = \alpha$:
\[
\phi(\alpha) < \phi(\alpha_0) + c_{1}(\alpha - \alpha_0) \cdot \phi^{\prime}(\alpha_0)
\]
Подставим значения $f$
\[
f(x - \alpha \nabla f(x)) <  f(x - 0 \cdot \nabla f(x)) - c_1 \alpha \| \nabla f(x) \|^2
\]
Пользуемся тем, что
\begin{multline*}
\phi^{\prime}(\alpha) =  f^{\prime}(x - \alpha\nabla f(x))(- \alpha \nabla f(x))^{\prime}_{\alpha} = \\  - f^{\prime}(x - \alpha\nabla f(x)) \nabla f(x) = - (\nabla f(x - \alpha \nabla f(x)))^T \nabla f(x)
\end{multline*}

Тогда $\phi^{\prime}(\alpha_0) = \phi^{\prime}(0) = - (\nabla f(x))^T \nabla f(x) = - \|\nabla f(x)\|^2$

Следовательно (при $\nabla f(x) \ne 0$) условие будет выполняться при $1 - \frac{L\alpha}{2} > c_{1} \Leftrightarrow \alpha < 2 \frac{1-c_1}{L}$ \\
Заметим, что предыдущий $\alpha$ мы отвергли, следовательно шаг в $\frac{1}{\beta}$ раз больший $\alpha$ не подходит;

Итого
\[\alpha \ge \min \left\{ \alpha_0; \frac{2\beta(1-c_1)}{L} \right\}\]

То есть $\alpha$ это шаг порядка $\frac{1}{L}$.
